\documentclass[12 pt, a4paper]{article}% тип документа, размер шрифта
\usepackage{cmap}	
\usepackage{hyperref}
\hypersetup{
    colorlinks=true,
    linkcolor=blue,
    urlcolor=blue,
}
\usepackage{mathtext}
\usepackage[T2A]{fontenc}%поддержка кириллицы в ЛаТеХ
\usepackage[utf8]{inputenc}%кодировка
\usepackage[english,russian]{babel}
\usepackage{indentfirst}

\usepackage{amsmath,amsfonts,amssymb,amsthm,mathtools} % AMS
\usepackage{amsmath}%удобная вёрстка многострочных формул, масштабирующийся текст в формулах, формулы в рамках и др.
\usepackage{amsfonts}%поддержка ажурного и готического шрифтов — например, для записи символа {\displaystyle \mathbb {R} } \mathbb {R} 
\usepackage{amssymb}%amsfonts + несколько сотен дополнительных математических символов
\frenchspacing%запрет длинного пробела после точки
\usepackage{setspace}%возможность установки межстрочного интервала
\usepackage{indentfirst}%пакет позволяет делать в первом абзаце после заголовка абзацный отступ
\onehalfspacing%установка полуторного интервала по умолчанию
\usepackage{graphicx}%подключение рисунков
\graphicspath{{images/}}%путь ко всем рисункам
\usepackage{caption}
\usepackage{float}%плавающие картинки
\usepackage{tikz} % это для чудо-миллиметровки
\usepackage{pgfplots}%для построения графиков
\pgfplotsset{compat=newest, y label style={rotate=-90},  width=10 cm}%версия пакета построения графиков, ширина графиков
\usepackage{pgfplotstable}%простое рисование табличек
\usepackage{lastpage}%пакет нумерации страниц
\usepackage{comment}%возможность вставлять большие комменты
\usepackage{python}%питончик
\usepackage{float}
%%%%% ПОЛЯ
\setlength\parindent{0pt} 
\usepackage[top = 2 cm, bottom = 2 cm, left = 1.5 cm, right = 1.5 cm]{geometry}
\setlength\parindent{0pt}
%%%%% КОЛОНТИТУЛЫ
\usepackage[shortlabels]{enumitem}

\usepackage{fancybox, fancyhdr}
\pagestyle{fancy} 
\fancyhead[L]{Легких Никита Б05-301}
\fancyhead[C]{2.2.6}
% ЛЮБАЯ ДОПОЛНИТЕЛЬНАЯ ИНФОРМАЦИЯ
\fancyfoot[R]{ФПМИ МФТИ}
\renewcommand{\footrulewidth}{0.3 mm}
 

\usepackage{graphicx}
\graphicspath{{pictures/}}
\DeclareGraphicsExtensions{.pdf,.png,.jpg}
\usepackage{wrapfig}
%%% Заголовок
%%% Новые команды
\newcommand{\z}[2]{{{\vspace{0.6cm} \textbf{{Задача {#1}}} \textit{{#2}\\ }}} }
\newcommand{\otv}{{\vspace{0.3cm} \textbf{Ответ: }}}
\newcommand{\uk}{\underline{\textit{Указание.}} }
\newcommand{\prim}{\underline{\textit{Примечание.}} }
\newcommand{\sol}[1]{{{\vspace{0.3cm} \textbf{{Задача {#1}} }\\ }}}



\begin{document} % конец преамбулы, начало документа

\begin{center}
\section*{Работа 2.2.6}
    \section*{Определение энергии активации по температурной зависимости вязкости жидкости}
\end{center}
\textbf{Цель работы:} 1) измерение скорости падения шариков при разной температуре жидкости; 2) вычисление вязкости жидкости по закону Стокса и расчет энергии активации. \\ \\
\textbf{В работе используются:} стеклянный цилиндр с исследуемой жидкостью (глицерин); термостат; секундомер; горизонтальный компаратор; микроскоп; мелкие шарики (диаметром около 1 мм).
\section*{Теоретическая справка}
Заметим некоторые особенности жидкости: 
\begin{itemize}
    \item жидкость обладает некоторыми свойствами газов и твердых тел: подобно газам, жидкости принимают форму сосуда, в котором они находятся, подобно твердым телам, они обладают сравнительно большой плотностью, с трудом поддаются сжатию.
    \item В жидкостях, как и в кристаллах, каждая молекула находится в потенциальной яме электрического поля, создаваемого окружающими молекулами
    \item  Глубина потенциальной ямы в жидкостях больше средней кинетической энергии колеблющейся молекулы, поэтому молекулы колеблются вокруг более или менее стабильных положений равновесия. Однако у жидкостей различие между этими двумя энергиями невелико, так что молекулы нередко выскакивают из «своей» потенциальной ямы и занимают место в другой.
    \item В жидкости присутствует ближний, но не дальний порядок.
\end{itemize}
Энергия актвации ($W$) - необходимая энергия для перехода молекулы из одного устойчивого состояния в другое. \\
Из распределени Больцмана: $n \sim e^{-W/kT} \Rightarrow \eta \sim A\cdot e^{W/kT}$ \\
\textbf{Формула Стокса} В нашем случае, очевидно, что сила сопротивления зависит от $\eta, r, \rho_{liq}, v$. Закон будет искать методов размерностей с помощью степенной формулы: $$F = A \eta^x r^y \rho_{liq}^z v^{\alpha}.$$ В общем виде система не имеет рещений, потому что уранвения только 3, а неизвестных 4, поэтому положим то что при наших измерениях скорости достаточно малы и будет считать $\alpha = 1$. Далее решаем систему уравнений: 
\begin{equation*}
 \begin{cases}
   1 = z + x
   \\
   1 = -x + 1 + y - 3z
   \\
   -2 = -x - 1
 \end{cases}
\end{equation*} и получаем: $x = 1 = y, z = 0 \Rightarrow$ $F = A \eta r v$, из строго решения получаем $A = 6\pi$. Итого: $F = 6\pi \eta r v$. \\
Нам также важно понимать, будет ли предположение о малой скорости верно, для этого надо оценивать значение числа Рейнольдса: $Re = \frac{\rho r v}{\eta}$, для того чтобы поток был достаточно ламинарным $Re < 10$.

Запишем второй закон Ньютона: $$Vg(\rho - \rho_{liq}) - 6\pi \eta r v = V\rho \frac{dv}{dt} \Rightarrow$$ $$v(t) = v(\infty) - [v(\infty) + v(0)]\cdot e^{-t / T}$$ 
$v(0)$ - начальная скорость \\
$v(\infty) = \frac{2}{9} g r^2 \frac{\rho - \rho_{liq}}{\eta}$ - установившаясь скорость \\ 
$\tau = \frac{V\rho}{6\pi \eta r} = \frac{2 r^2 \rho}{9 \eta}$ \\ 
$S = v(\infty) \tau \cdot (\frac{t}{\tau} - 1 + e^{-t/T})$ \\
Если $t > 5\tau$, то процесс можно считать установившимся.
\section*{Задание}
Значения установки: плотность стали - $\rho_1 = 7800$ кг/м$^3$, плотность стекла - $\rho_2 = 2600$ кг/м$^3$.
$L = (10\pm0.3)$ см \\
$D = (4\pm0.5)$ см \\
$\rho_{liq}(T) = \frac{571}{450} - 4.(4)\cdot10^{-4} \cdot T$ \\ 
$\eta = \frac{2}{9} g r^2 \frac{\rho - \rho_{liq}}{v(\infty)} \approx \frac{2}{9} g r^2 \frac{\rho - \rho_{liq}}{L} t$ \\
\prim  Шарики под номерами 1 и 2 - всегда стеклянные, а под 3 и 4 - всегда железные
\begin{enumerate}
    \item $T = 24.41 ^oC \rightarrow \rho_{liq} = 1.258$ г/см$^3$
    \begin{table}[h!] 
   \begin{center}
    \begin{tabular}{|*{11}{l|}} \hline
     $N$ & $d_1$, мм  & $d_2$, мм & $d_3$, мм & $<d>$, мм & $t_1$, c & $t_2$, c & $\eta_1$, Па$\cdot$с & $\eta_2$, Па$\cdot$с & $\eta$, Па$\cdot$с & $\delta\eta$, Па$\cdot$с \\ \hline
     1 & 2.02 & 2.04 & 2.00 & 2.02 & 23.8 & 25.8 & 77.3 & 78.7 & 78 & 0.5 \\ \hline
     2 & 2.14 & 2.18 & 2.13 & 2.15 & 23.3 & 24.9 & 75.1 & 77.0 & 76 & 1.0\\ \hline
     3 & 0.64 & 0.62 & 0.63 & 0.63 & 27.8 & 29.2 & 48.7 & 49.3 & 49 & 0.3\\
     \hline
     4 & 0.75 & 0.73 & 0.74 & 0.74 & 25.8 & 27.7 & 46.0 & 44.0 & 45 & 1.0\\
     \hline
    \end{tabular}
   \end{center}
  \end{table} \\
  $Re_1 = \frac{\rho_{liq} r_1 v}{\eta} = \frac{\rho_{liq} r_1 L}{\eta t} \approx 1.6\cdot10^{-7} $,
  $Re_2 = \frac{\rho_{liq} r_2 v}{\eta} = \frac{\rho_{liq} r_2 L}{\eta t} \approx 8\cdot10^{-7} \Rightarrow$ течение ламинарное
  $\tau = \frac{2 <r>^2 \rho_{liq}}{9 \eta} \approx 2$ нс $\Rightarrow$ можно считать что движение почти сразу установившееся
  \item $T = 30.80 ^oC \rightarrow \rho_{liq} = 1.255$ г/см$^3$
    \begin{table}[h!] 
   \begin{center}
    \begin{tabular}{|*{11}{l|}} \hline
     $N$ & $d_1$, мм  & $d_2$, мм & $d_3$, мм & $<d>$, мм & $t_1$, c & $t_2$, c & $\eta_1$, Па$\cdot$с & $\eta_2$, Па$\cdot$с & $\eta$, Па$\cdot$с & $\delta\eta$, Па$\cdot$с \\ \hline
     1 & 2.10 & 2.12 & 2.12 & 2.11 & 14.7 & 15.4 & 47.4 & 48.7 & 48 & 0.5 \\ \hline
     2 & 2.10 & 2.10 & 2.10 & 2.10 & 15.3 & 16.2 & 50.8 & 49.2 & 50 & 0.8\\ \hline
     3 & 0.70 & 0.80 & 0.74 & 0.75 & 14.1 & 14.7 & 24.3 & 25.7 & 25 & 0.6\\
     \hline
     4 & 0.74 & 0.76 & 0.75 & 0.75 & 19.4 & 19.3 & 33.0 & 35.3 & 34 & 1.1\\
     \hline
    \end{tabular}
   \end{center}
  \end{table} \\
  $Re_1 = \frac{\rho_{liq} r_1 v}{\eta} = \frac{\rho_{liq} r_1 L}{\eta t} \approx 2.3\cdot10^{-7} $,
  $Re_2 = \frac{\rho_{liq} r_2 v}{\eta} = \frac{\rho_{liq} r_2 L}{\eta t} \approx 5\cdot10^{-7} \Rightarrow$ течение ламинарное
  $\tau = \frac{2 <r>^2 \rho_{liq}}{9 <\eta>} \approx 2$ нс $\Rightarrow$ можно считать что движение почти сразу установившееся
  \item $T = 35.04 ^oC \rightarrow \rho_{liq} = 1.253$ г/см$^3$
    \begin{table}[h!] 
   \begin{center}
    \begin{tabular}{|*{11}{l|}} \hline
     $N$ & $d_1$, мм  & $d_2$, мм & $d_3$, мм & $<d>$, мм & $t_1$, c & $t_2$, c & $\eta_1$, Па$\cdot$с & $\eta_2$, Па$\cdot$с & $\eta$, Па$\cdot$с & $\delta\eta$, Па$\cdot$с \\ \hline
     1 & 2.04 & 2.06 & 2.04 & 2.06 & 12.3 & 12.0 & 39.5 & 40.5 & 40 & 0.5 \\ \hline
     2 & 2.10 & 2.08 & 2.10 & 2.09 & 11.1 & 11.9 & 35.2 & 36.9 & 36 & 0.8\\ \hline
     3 & 0.72 & 0.70 & 0.72 & 0.72 & 22.8 & 25.0 & 30.3 & 25.1 & 27 & 2.2\\
     \hline
     4 & 0.70 & 0.72 & 0.74 & 0.72 & 15.8 & 16.3 & 33.0 & 35.3 & 28 & 0.3\\
     \hline
    \end{tabular}
   \end{center}
  \end{table} \\
  $Re_1 = \frac{\rho_{liq} r_1 v}{\eta} = \frac{\rho_{liq} r_1 L}{\eta t} \approx 5.6\cdot10^{-7} $,
  $Re_2 = \frac{\rho_{liq} r_2 v}{\eta} = \frac{\rho_{liq} r_2 L}{\eta t} \approx 6.3\cdot10^{-7} \Rightarrow$ течение ламинарное
  $\tau = \frac{2 <r>^2 \rho_{liq}}{9 <\eta>} \approx 1$ нс $\Rightarrow$ можно считать что движение почти сразу установившееся
  \item $T = 39.75 ^oC \rightarrow \rho_{liq} = 1.251$ г/см$^3$
    \begin{table}[h!] 
   \begin{center}
    \begin{tabular}{|*{11}{l|}} \hline
     $N$ & $d_1$, мм  & $d_2$, мм & $d_3$, мм & $<d>$, мм & $t_1$, c & $t_2$, c & $\eta_1$, Па$\cdot$с & $\eta_2$, Па$\cdot$с & $\eta$, Па$\cdot$с & $\delta\eta$, Па$\cdot$с \\ \hline
     1 & 2.00 & 2.08 & 2.08 & 2.05 & 8.58 & 9.07 & 27.5 & 28.7 & 28 & 0.5 \\ \hline
     2 & 2.10 & 2.10 & 2.10 & 2.10 & 8.87 & 9.01 & 29.0 & 29.3 & 29 & 0.2\\ \hline
     3 & 0.76 & 0.76 & 0.76 & 0.76 & 8.57 & 9.70 & 13.8 & 16.4 & 15 & 1.5\\
     \hline
     4 & 0.70 & 0.72 & 0.74 & 0.72 & 12.10 & 12.67 & 19.7 & 20.3 & 20 & 0.3\\
     \hline
    \end{tabular}
   \end{center}
  \end{table} \\
  $Re_1 = \frac{\rho_{liq} r_1 v}{\eta} = \frac{\rho_{liq} r_1 L}{\eta t} \approx 7\cdot10^{-7} $,
  $Re_2 = \frac{\rho_{liq} r_2 v}{\eta} = \frac{\rho_{liq} r_2 L}{\eta t} \approx 9.1\cdot10^{-7} \Rightarrow$ течение ламинарное
  $\tau = \frac{2 <r>^2 \rho_{liq}}{9 <\eta>} \approx 1$ нс $\Rightarrow$ можно считать что движение почти сразу установившееся
  \item $T = 44.67 ^oC \rightarrow \rho_{liq} = 1.249$ г/см$^3$
    \begin{table}[h!] 
   \begin{center}
    \begin{tabular}{|*{11}{l|}} \hline
     $N$ & $d_1$, мм  & $d_2$, мм & $d_3$, мм & $<d>$, мм & $t_1$, c & $t_2$, c & $\eta_1$, Па$\cdot$с & $\eta_2$, Па$\cdot$с & $\eta$, Па$\cdot$с & $\delta\eta$, Па$\cdot$с \\ \hline
     1 & 2.10 & 2.12 & 2.12 & 2.12 & 6.33 & 5.75 & 27.5 & 28.7 & 20 & 0.5 \\ \hline
     2 & 2.10 & 2.10 & 2.10 & 2.10 & 5.88 & 6.05 & 29.0 & 29.3 & 19 & 0.2\\ \hline
     3 & 0.70 & 0.80 & 0.74 & 0.75 & 8.24 & 8.34 & 14.8 & 15.4 & 14 & 0.3\\
     \hline
     4 & 0.74 & 0.76 & 0.76 & 0.76 & 7.87 & 7.52 & 19.7 & 20.3 & 13 & 0.3\\
     \hline
    \end{tabular}
   \end{center}
  \end{table} \\
  $Re_1 = \frac{\rho_{liq} r_1 v}{\eta} = \frac{\rho_{liq} r_1 L}{\eta t} \approx 7\cdot10^{-7} $,
  $Re_2 = \frac{\rho_{liq} r_2 v}{\eta} = \frac{\rho_{liq} r_2 L}{\eta t} \approx 9.1\cdot10^{-7} \Rightarrow$ течение ламинарное
  $\tau = \frac{2 <r>^2 \rho_{liq}}{9 <\eta>} \approx 1$ нс $\Rightarrow$ можно считать что движение почти сразу установившееся
\end{enumerate}

\textbf{Линеаризация и график} \\
Вспоминаем формулу $\eta \sim A\cdot e^{W/kT}$, для нее очевидно линеаризация для нахождения $W \rightarrow$ \\ $\ln{\eta}$ от $\frac{1}{T}$, построим график: 
\begin{figure}[h]
\centering
\includegraphics[width=0.7\linewidth]{2.2.6 Определение энергии активации по температурной зависимости вязкости жидкости/график 1.png}
\label{fig:mpr}
\end{figure}

Из графика получаем значения: $$k_{стекла} = \frac{W}{k} = 5.8 \; К\cdot10^3$$
$$k_{сталь} = \frac{W}{k} = 6.5 \; К\cdot10^3$$
Погрешность же оценим в помощью МНК включив в нее приборную погрещность, потому что кресты погрешностей слишком малы более точную оценку даст именно такой метод.\\
Стальные шарики: 
$$\Delta k_{МНК} = \sqrt{\frac{<\ln(\eta)^2>-<\ln(\eta)>^2}{<1/T^2>-<1/T>^2}-k^2}/\sqrt{N} = 0.4 \; К\cdot10^3$$ $$\Delta k_{стат} = \frac{<\delta \eta>}{<\eta>}\cdot <k> = 0.3\; К\cdot10^3$$
$$\Delta k_{сталь} = \sqrt{\Delta k_{МНК}^2 + \Delta k_{стат}^2} = 0.5 \; К\cdot10^3$$
Стеклянные шарики: 
$$\Delta k_{МНК} = \sqrt{\frac{<\ln(\eta)^2>-<\ln(\eta)>^2}{<1/T^2>-<1/T>^2}-k^2}/\sqrt{N} = 0.4 \; К\cdot10^3$$ $$\Delta k_{стат} = \frac{<\delta \eta>}{<\eta>}\cdot <k> = 0.2\; К\cdot10^3$$
$$\Delta k_{стекло} = \sqrt{\Delta k_{МНК}^2 + \Delta k_{стат}^2} = 0.4 \; К\cdot10^3$$
Итого: $$k_{сталь} = 6.5\; К\cdot10^3 \Rightarrow W_{сталь} = k_{сталь} \cdot k = 0.55\; эВ$$
$$\Delta k_{сталь} = 0.5\; К\cdot10^3 \Rightarrow \Delta W_{сталь} = W_{сталь} \cdot \epsilon_{k_{сталь}} =  0.04 \;эВ$$
$$k_{стекло} = 5.8\; К\cdot10^3 \Rightarrow W_{стекло} = k_{стекло} \cdot k = 0.50 \;эВ$$
$$\Delta k_{стекло} = 0.4\; К\cdot10^3 \Rightarrow \Delta W_{стекло} = W_{стекло} \cdot \epsilon_{k_{стекло}} =  0.03 \;эВ$$
\section*{Вывод}
Мы получили ответы: $W_{сталь} = 0.55\pm0.04 \; эВ$ и $W_{стекло} = 0.5\pm0.03 \; эВ$, в ворота каждого из которых попадает значение взятое из справочных материалов: $W_{справ} = 0.52 \; эВ$, и которые между собой отличаются лишь на $(\frac{W_{сталь}}{W_{стекло}} - 1) / 100 = 10\%$.
\end{document}